\clearpage
\pdfbookmark[0]{Resumo}{resumo}
\chapter*{RESUMO}
\addcontentsline{toc}{chapter}{RESUMO}

Este trabalho compara estruturas de cooperação entre canais de informação que emergem após o treinamento de classificadores com amostras de duas origens. Na condição de referência, os classificadores são treinados com amostras reais; nas condições sintéticas, os mesmos classificadores são treinados com amostras produzidas por modelos probabilísticos parametrizados a partir dos dados reais de treinamento. Em cada condição, a estrutura de cooperação é caracterizada pela variação do ranking dos melhores subconjuntos de canais à medida que aumenta a cardinalidade do conjunto de treinamento. Investiga-se, então, em que medida a estrutura obtida nas condições sintéticas reproduz aquela obtida na condição real. O protocolo compara três condições de treinamento---dados reais, Gaussiana única condicional à classe e mistura de gaussianas---avaliadas sobre o mesmo conjunto de teste real fixo. Todos os subconjuntos não vazios de canais são examinados ao longo de uma grade de tamanhos amostrais por classe por meio de simulação de Monte Carlo com amostras incrementalmente aninhadas. Os experimentos utilizam os datasets Occupancy Detection, UCI Air Quality e SUPPORT2, com SVM linear, regressão logística, Gaussian Naive Bayes e, adicionalmente no SUPPORT2, Random Forest. A preservação é avaliada por correlação de ranking, \WinnersAgreement{} entre comparações par a par e similaridade progressiva das matrizes de $N^{\star}$, definidas pelo último cruzamento observado entre curvas de perda. Os resultados mostram que a fidelidade estrutural é multidimensional: ranking e relações de vencedores podem ser fortemente preservados mesmo quando a similaridade de $N^{\star}$ é moderada ou baixa. No Air Quality, o Gaussian Naive Bayes apresentou preservação praticamente integral do ranking e dos vencedores; no SUPPORT2, o Random Forest combinou baixo desempenho preditivo absoluto com a maior similaridade de $N^{\star}$ entre os classificadores daquele cenário, apesar de milhares de cruzamentos definidos. O treinamento com amostras da Gaussiana única foi competitivo ou superior ao treinamento com amostras do GMM em algumas métricas, embora o GMM reproduzisse visualmente melhor certas geometrias. Conclui-se que desempenho preditivo, fidelidade geométrica e preservação da cooperação são propriedades distintas. O trabalho não demonstra redução de custos, mas fornece objetos estruturais para futura análise custo-eficiência e para investigação de projeções em regimes amostrais maiores.

\vspace{\onelineskip}
\noindent\textbf{Palavras-chave}: canais de informação; classificação cooperativa; dados sintéticos; seleção de subconjuntos; simulação de Monte Carlo; fidelidade estrutural.
